\section{Robustez del verificador}
\label{sec:robustez}

\pendiente{describir qué valida el verificador, cómo interpreta la
respuesta y cómo implementa la interfaz
\texttt{reward(prompt, completion) -> float}.}

\subsection{Equivalencias y pruebas adversariales}
\guia{Comprobar que se acepten respuestas correctas equivalentes y
se rechacen respuestas incorrectas. El enunciado valora casos
diseñados para engañar al verificador, incluyendo respuestas
contradictorias, explotación del formato y errores del parser.
La tabla propone categorías de prueba, no resultados ya obtenidos.}

\begin{table}[htbp]
  \centering\small
  \caption{Pruebas de robustez del verificador.}
  \label{tab:robustez}
  \begin{tabularx}{\linewidth}{@{}YYYY@{}}
    \toprule
    Caso probado & Comportamiento esperado & Resultado observado & Evidencia \\
    \midrule
    Respuesta correcta equivalente & Aceptar & \dato & \dato \\
    Respuesta incorrecta & Rechazar & \dato & \dato \\
    Respuestas contradictorias & Rechazar & \dato & \dato \\
    Formato inválido o ambiguo & Definir política & \dato & \dato \\
    Intento de engaño al parser & Rechazar el engaño & \dato & \dato \\
    \bottomrule
  \end{tabularx}
\end{table}

\subsection{Vulnerabilidades y limitaciones}
\pendiente{reportar el número de casos evaluados, falsos positivos
(respuestas incorrectas aceptadas), falsos negativos (respuestas
correctas rechazadas), problemas encontrados y correcciones.}

\guia{Distinguir los fallos del modelo de los del verificador.
Si se modifica este último, indicar si se recalcularon las métricas
afectadas. Declarar qué ataques o equivalencias aún no se han probado.}
